\documentclass[UTF8,a4paper,11pt]{ctexart}
\usepackage[margin=22mm]{geometry}
\usepackage{amsmath,amssymb,booktabs,longtable,hyperref}
\setlength{\parindent}{0pt}
\setlength{\parskip}{6pt}
\begin{document}
\begin{center}\LARGE\bfseries 当前 5\(\times\)5 非平面拼接的刚性证明\end{center}
\begin{center}完整构造文件说明 · 非刚性反例 · 可精确复核的数学证明\end{center}
\begin{center}2026 年 10 月 16 日\end{center}
\section*{1  结论与适用范围}
本文只讨论项目中已确认的那一个 5\(\times\)5 初始空间拼接：顶点坐标见第 4 节，25 个空间四边形的连接关系固定，每个四边形整体保持刚性，相邻面沿整条公共边连接。

定理：在上述条件下，从指定初态出发不存在非平凡的连续刚性折叠。整体的平移、转动不算折叠。结论不依赖 F 的原圆，也不依赖具体驱动点、转动速度或动画算法。

这不是“任意非平面 5\(\times\)5 都不可折叠”的定理。改变初始坐标、面片朝向、连接关系，或者允许一个空间四边形内部的两个三角形相对转动，均属于另一个问题。本文也不排除与初态不能连续相连的其他离散拼装。

\section*{2  交付的完整构造文件}
5x5 nonplanar folding full construction.ggb 是现有联动文件的完整代数区版本，不是新实现的刚性折叠。文件保留全部 203 条构造定义，其中原始 2\(\times\)2 的 48 条定义逐条保留；本次没有修改点的位置关系、圆、交点或联动公式。

\begin{longtable}{lll}	oprule
对象 & 数量 & 代数区处理 \\ \midrule
点（含中间构造点） & 66 & 全部显示定义和数值 \\
线段（含隐藏的面内边） & 247 & 全部显示端点及定义 \\
多边形 / 三角面 & 59 & 保留原有及显示用面片 \\
圆 / 圆锥曲线 & 8 & 保留原圆及辅助圆 \\
平面；向量 & 3；2 & 全部保留 \\
合计 & 385 & 均不再作为隐藏的辅助对象 \\
\bottomrule\end{longtable}
左侧按构造顺序显示“定义与数值”。三维画面仍隐藏多数辅助点线圆，避免遮挡；可以在代数区逐个开启其可见性。原有 q1 = Polygon(A,B,C,D) 因四点不共面可能显示“?”：其定义仍被原样保留，实际空间面由两片同色三角形显示。不要把这个旧多边形的无定义数值误认为构造丢失。

文件中的驱动点是 F = Point(e)。P = Intersect(r\textsubscript{1},t) 是依赖交点，不是独立驱动点。现有联动的变形问题仍保留，便于检查和复现。

\newpage
\section*{3  当前联动为什么不刚性}
外围 27 个网格点不是把整个刚性面统一作同一个等距变换而得到的，而是分别对不同的中点作点对称。单次点对称是等距变换，但不同顶点采用不同中心，不能推出拼成的新四边形保持全等。

例如第一行第二列黄色面按顺序取顶点 W\textsubscript{01}、W\textsubscript{02}、D、O。文件中的中点与对称命令给出以下精确关系：


\[
W_{01}=O+F-K,\qquad W_{02}=D+A-G=D-G.
\]

初态 F = (-1,1,\(\sqrt{2}\)) 时，这个面的四边均长 2，短对角线 OW\textsubscript{02} 长 2。将 F 沿原圆移到以下位置，并取与初态连续对应的原交点分支，原始构造给出：


\[
F=(0,1,\sqrt{3}),\qquad G=(1,-3/2,\sqrt{3}/2),
\]


\[
O=(-2\sqrt{6}/3,\,2,\,2\sqrt{3}/3),\qquad W_{02}=(-1,7/2,-\sqrt{3}/2).
\]

这些坐标直接代入原圆及平面交点条件即可核验。于是短对角线满足：


\[
|W_{02}-O|^2=(-1+2\sqrt{6}/3)^2+(3/2)^2+(-7\sqrt{3}/6)^2
\]


\[
=10-\frac{4\sqrt{6}}{3}\ne 4,\qquad |W_{02}-O|\approx 2.594998.
\]

这已经是一个精确的非刚性反例，无须依赖图像观感。隐藏面内对角线不会消除这个点间距离的变化。

\section*{原生 GeoGebra 的另一组复核}
以圆参数 \(\theta\) 定义 F(\(\theta\)) = (\(\sqrt{3}\) cos \(\theta\),1,\(\sqrt{3}\) sin \(\theta\))，初态 \(\theta\)\textsubscript{0} = arccos(-1/\(\sqrt{3}\))。正向 10\({}^\circ\) 指 \(\theta\)\textsubscript{0} + 10\({}^\circ\)，不是任意投影视角下的顺时针。

\begin{longtable}{lll}	oprule
同一黄色面的长度 & 初态 & 圆参数增加 10\({}^\circ\) \\ \midrule
边 W\textsubscript{01}W\textsubscript{02} & 2.000000000 & 2.029394026 \\
长对角线 W\textsubscript{01}D & 2.828427125 & 3.026204354 \\
短对角线 W\textsubscript{02}O & 2.000000000 & 1.720618297 \\
\bottomrule\end{longtable}
对初态两侧 \(\pm\)25\({}^\circ\) 的 101 个位置逐一读取 GeoGebra 坐标：原始四面的六个距离保持不变（浮点误差约 1.8\(\times\)10\textsuperscript{-15}）；新增 21 面均出现距离变化。该采样仅作为现有文件的诊断，不作为下面不可折叠定理的证明。

\newpage
\section*{4  精确模型与刚性约束}
全局顶点写作 p(r,c)，其中 r,c = 0,…,5；单元面为四点 p(r,c)、p(r,c+1)、p(r+1,c+1)、p(r+1,c)，其中 r,c = 0,…,4。相邻面共用同一顶点，因而不允许公共边的两端被复制成不同位置。

为避免根式运算，使用缩放坐标 x(r,c)。物理坐标、内积和精确初态如下；所有矩阵运算都在有理数域进行。


\[
p_{rc}=Sx_{rc},\qquad S=\mathrm{diag}(1,1,\sqrt{2}),\qquad M=S^TS=\mathrm{diag}(1,1,2).
\]


\[
\|x\|_M^2=x^TMx,\qquad a=(1,-3,1),\quad b=(3,-1,-1).
\]


\[
x^0_{rc}=b_{\epsilon_r\epsilon_c}+\lfloor(r-1)/2\rfloor a+\lfloor(c-1)/2\rfloor b,
\]


\[
\epsilon_r=(r-1)\;\mathrm{mod}\;2,\qquad \epsilon_c=(c-1)\;\mathrm{mod}\;2,
\]


\[
b_{00}=(-2,2,0),\ b_{01}=(0,2,0),\ b_{10}=(-1,1,1),\ b_{11}=(0,0,0).
\]

这里的整数坐标与 5x5 nonplanar static.ggb 的根式定义逐点一致；不是从屏幕图像估计坐标，也不是把浮点测量误差当作精确几何。完整构造版在上述 F 初态计算得到同一拼接。

每一个空间四边形整体刚性，要求四个顶点的全部六个点间距离恒定，即四条边和两条对角线均不变。对当前每个单元，六个距离平方的集合为 \{4,4,4,4,4,8\}，且四点不共面。


\[
(x_i(t)-x_j(t))^TM(x_i(t)-x_j(t))=(x_i^0-x_j^0)^TM(x_i^0-x_j^0).
\]

上式对每个单元内的每一对顶点成立。把这些约束记作“杆”只是数学表示，不是在 GGB 中添加可见线段，也不允许面内出现新折痕。25 个单元合计形成 110 个不同的点对约束。

只要证明这些必要距离约束已禁止非平凡连续运动，就足以禁止真实刚性面片折叠；是否还要求无自交，并不影响这个否定结论。

\newpage
\section*{5  九面子结构与参考面的固定}
先选取全局 r,c = 1,…,4 的 16 个顶点，即第 2 至第 4 行、列的九个单元。将这 16 个顶点按行优先编号为 1,…,16；下表给出缩放坐标。实际高度为表中第三坐标乘 \(\sqrt{2}\)。

\begin{longtable}{lll}	oprule
编号 & 完整构造版名称 & 缩放坐标 \\ \midrule
1 & O & (-2, 2, 0) \\
2 & D & (0, 2, 0) \\
3 & C & (1, 1, -1) \\
4 & W\_\{14\} & (3, 1, -1) \\
5 & F & (-1, 1, 1) \\
6 & A & (0, 0, 0) \\
7 & B & (2, 0, 0) \\
8 & W\_\{24\} & (3, -1, -1) \\
9 & K & (-1, -1, 1) \\
10 & G & (1, -1, 1) \\
11 & Q & (2, -2, 0) \\
12 & W\_\{34\} & (4, -2, 0) \\
13 & W\_\{41\} & (0, -2, 2) \\
14 & W\_\{42\} & (1, -3, 1) \\
15 & W\_\{43\} & (3, -3, 1) \\
16 & W\_\{44\} & (4, -4, 0) \\
\bottomrule\end{longtable}
固定参考面 D,C,A,B，即编号 2、3、6、7 的四个顶点。该面自身刚性且四点不共面，所以任意连续刚性运动都可通过一个随时间变化的整体刚体变换，把这四点固定在初态位置。此操作只去掉整体平移、转动，不额外限制折叠。

固定后余下 12 点、36 个坐标自由量。该九面子结构共有 42 个不同的点对约束，详列于附录。接下来证明固定参考面后，其初态是距离方程解集中的孤立点。

\newpage
\section*{6  秩与唯一无穷小方向的精确证书}
记固定后的位移列向量为 h\(\in\)R\textsuperscript{36}；其分量按顶点编号递增、每点 x,y,z 顺序排列，跳过固定点 2、3、6、7。对点对 i<j，定义矩阵 R 的对应行为以下线性型，固定点的位移视为零。


\[
(Rh)_{ij}=(x_i^0-x_j^0)^TM(h_i-h_j),\qquad R\in\mathbb{Q}^{42\times36}.
\]

行按点对 (i,j) 的字典序排列。取下表 v 为完整的 16 点缩放速度；固定点的零向量删除后得到 \(\bar v\)\(\in\)R\textsuperscript{36}。直接代入给出 R\(\bar v\) = 0。

\begin{longtable}{llll}	oprule
点号 & v\textsubscript{i} & 点号 & v\textsubscript{i} \\ \midrule
1 & (0, 0, -2) & 9 & (0, 0, 0) \\
2 & (0, 0, 0) & 10 & (0, 2, 1) \\
3 & (0, 0, 0) & 11 & (0, 0, 2) \\
4 & (0, -2, -1) & 12 & (0, 0, 0) \\
5 & (-2, 0, -1) & 13 & (0, 4, 2) \\
6 & (0, 0, 0) & 14 & (2, 2, 4) \\
7 & (0, 0, 0) & 15 & (2, 4, 3) \\
8 & (-2, -2, 0) & 16 & (0, 0, 4) \\
\bottomrule\end{longtable}
为使秩的结论能够完全复核，给出一个显式非零子式：从 R 删除第 4、6、7、9、10、17、41 行，再删除第 36 列，得到 35\(\times\)35 矩阵 R*。整数消元的精确结果是：


\[
\det R_*=-68719476736=-2^{36}\ne0.
\]

故 rank R\(\ge\)35；另一方面 \(\bar v\)\(\ne\)0 且 R\(\bar v\)=0，故 rank R\(\le\)35。因此：


\[
\mathrm{rank}\,R=35,\qquad\ker R=\mathrm{span}\{\bar v\}.
\]

这一步不能单独推出存在可折叠运动。“存在一个无穷小方向”只说明距离约束的一阶项可以抵消；接下来检查这个方向能否对应真正相邻构型。

\newpage
\section*{7  二次项证书}
按下表给点对赋权 \(\omega\)\textsubscript{ij}，其余点对权重取零。这些正负整数是代数检验权重，不表示必须在真实折纸上施加外力。

\begin{longtable}{llll}	oprule
点对 & 权重 & 点对 & 权重 \\ \midrule
(2,5) & 1 & (8,12) & 1 \\
(3,8) & 1 & (9,10) & -1 \\
(5,6) & -1 & (9,14) & 1 \\
(5,9) & 1 & (10,11) & 1 \\
(6,10) & 1 & (10,14) & -1 \\
(6,11) & -2 & (11,12) & -1 \\
(7,8) & -1 & (11,15) & -1 \\
(7,10) & -2 & (12,15) & 1 \\
(7,11) & 1 & (14,15) & 1 \\
\bottomrule\end{longtable}
逐列代入第 6 节矩阵，得到以下精确恒等式。这里 \(\omega\) 对固定后的矩阵成立；固定点上的反作用无需另作约束。


\[
\omega^T R=0.
\]

对任意位移 h 定义二次型 Q，仍令四个固定点的位移为零。把第 6 节的整数速度逐项代入，得到：


\[
Q(h)=\sum_{(i,j)}\omega_{ij}(h_i-h_j)^TM(h_i-h_j),
\]


\[
Q(\bar v)=32,\qquad Q(t\bar v)=32t^2.
\]

Q 未必在整个 36 维空间半正定；本证明只需它在 ker R 的所有非零向量上严格为正。以上一维核与整数值 32 已经充分。不能把 Q 的这一性质错误推广成全空间正定。

附带的核验程序不使用浮点 SVD 判秩：R、\(\bar v\)、\(\omega\)、非零子式和数值 32 全部采用 SymPy 有理数 / 整数运算复核。

\newpage
\section*{8  严格证明九面子结构局部锁死}
完整距离方程在初态附近具有精确的二次展开，而非截断近似。对每个点对，令 q\textsubscript{ij}(h) 为位移差的 M-范数平方，则：


\[
2Rh+q(h)=0,\qquad q_{ij}(h)=(h_i-h_j)^TM(h_i-h_j).
\]

左乘 \(\omega\)\textsuperscript{T}，由 \(\omega\)\textsuperscript{T}R=0 可知，任意满足全部距离方程的位移都必须满足 Q(h)=0。

反证：若固定参考面后，初态不是解集中的孤立点，则存在非零精确解 h\textsuperscript{(n)}\(\to\)0。取 \(\varepsilon\)\textsubscript{n}=\(\Vert\)h\textsuperscript{(n)}\(\Vert\)\textsubscript{2}，并令 u\textsuperscript{(n)}=h\textsuperscript{(n)}/\(\varepsilon\)\textsubscript{n}。单位球紧致，故可取子列使 u\textsuperscript{(n)}\(\to\)u 且 \(\Vert\)u\(\Vert\)\textsubscript{2}=1。


\[
2Ru^{(n)}+\varepsilon_n q(u^{(n)})=0,\qquad Q(u^{(n)})=0.
\]

令 n\(\to\)\(\infty\)，连续性给出 Ru=0、Q(u)=0。由 ker R=span\{\(\bar v\)\}，有 u=t \(\bar v\)；由于 \(\Vert\)u\(\Vert\)\textsubscript{2}=1，必有 t\(\ne\)0。然而：


\[
0=Q(u)=Q(t\bar v)=32t^2>0,
\]

矛盾。因此存在一个初态邻域，在该邻域内固定参考面的距离方程只有初态这一解。九面子结构局部刚性，结论得证。

\section*{为什么这已排除“换速度、换参数再折”}
上述证明直接讨论相邻构型的存在性，没有假设运动可微、初速度非零或二阶导数存在。因此，初速度取零、极慢启动、换参数、改变驱动点，都不能绕过结论。固定参考面后不存在任何可供连续运动离开的相邻解。

这也解释了一阶线性化为何会误导：R\(\bar v\)=0 的确成立，但二次距离项在该方向上无法同时抵消。本文采用的是精确解集的孤立性证明，不是把一阶“看起来能动”当作可折叠性。

该论证是应力 / 二阶刚性思想在本实例中的直接实现。Connelly-Whiteley [1] 的 Corollary 5.2.2 给出杆框架二阶应力检验，Theorem 4.3.1 给出二阶刚性蕴含刚性。这里已给出独立的邻域证明，不需仅凭文献定理或数值优化失败作结论。

\newpage
\section*{9  从九个面推广到完整 5\(\times\)5}
九面子结构固定后，全局 r,c=1,…,4 的 16 个顶点固定。下面按表格顺序固定其余 20 点。每个目标点到所列三个中心的距离，均是某个刚性面中的边长或对角线长度；三中心在初态不共线，并在本步骤开始前已经固定。

\begin{longtable}{lll}	oprule
步骤 & 目标 (r,c) & 三个已固定中心的 (r,c) \\ \midrule
1 & (0, 2) & (1, 1)；(1, 2)；(1, 3) \\
2 & (0, 3) & (1, 2)；(1, 3)；(1, 4) \\
3 & (5, 2) & (4, 1)；(4, 2)；(4, 3) \\
4 & (5, 3) & (4, 2)；(4, 3)；(4, 4) \\
5 & (2, 0) & (1, 1)；(2, 1)；(3, 1) \\
6 & (3, 0) & (2, 1)；(3, 1)；(4, 1) \\
7 & (2, 5) & (1, 4)；(2, 4)；(3, 4) \\
8 & (3, 5) & (2, 4)；(3, 4)；(4, 4) \\
9 & (0, 1) & (1, 1)；(1, 2)；(0, 2) \\
10 & (0, 4) & (1, 3)；(1, 4)；(0, 3) \\
11 & (5, 1) & (4, 1)；(4, 2)；(5, 2) \\
12 & (5, 4) & (4, 3)；(4, 4)；(5, 3) \\
13 & (1, 0) & (1, 1)；(2, 1)；(2, 0) \\
14 & (4, 0) & (3, 1)；(4, 1)；(3, 0) \\
15 & (1, 5) & (1, 4)；(2, 4)；(2, 5) \\
16 & (4, 5) & (3, 4)；(4, 4)；(3, 5) \\
17 & (0, 0) & (0, 1)；(1, 0)；(1, 1) \\
18 & (0, 5) & (0, 4)；(1, 4)；(1, 5) \\
19 & (5, 0) & (4, 0)；(4, 1)；(5, 1) \\
20 & (5, 5) & (4, 4)；(4, 5)；(5, 4) \\
\bottomrule\end{longtable}
引理：到三个固定且不共线中心的距离均固定时，点只能处于至多两个位置。证明：三条球面方程两两相减，得到法向量线性独立的两张平面；其交集是一条直线，该直线与第一个球面至多相交两点。故原位置在解集中孤立，连续运动不能跳到另一交点。

逐项应用该引理，36 个顶点全部局部固定。每一步的三中心非共线性、点对确属某个单元的六距离约束，均在核验程序中精确检查。因此整个 5\(\times\)5 在去掉整体刚体运动后局部孤立。由连续性，不存在从该初态出发的非平凡连续刚性折叠。证毕。

\newpage
\section*{附录 A  矩阵的完整行顺序}
以下为第 6 节 42 个点对的字典序。每行对应 (x\textsubscript{i}\textsuperscript{0}-x\textsubscript{j}\textsuperscript{0})\textsuperscript{T}M(h\textsubscript{i}-h\textsubscript{j})。结合坐标表、固定点集合和列顺序，R 的每一个整数元素均已唯一确定；无需另行下载一个大矩阵。

\begin{longtable}{llllll}	oprule
行号 & 点对 & 行号 & 点对 & 行号 & 点对 \\ \midrule
1 & (1, 2) & 15 & (5, 9) & 29 & (9, 14) \\
2 & (1, 5) & 16 & (5, 10) & 30 & (10, 11) \\
3 & (1, 6) & 17 & (6, 7) & 31 & (10, 13) \\
4 & (2, 3) & 18 & (6, 9) & 32 & (10, 14) \\
5 & (2, 5) & 19 & (6, 10) & 33 & (10, 15) \\
6 & (2, 6) & 20 & (6, 11) & 34 & (11, 12) \\
7 & (2, 7) & 21 & (7, 8) & 35 & (11, 14) \\
8 & (3, 4) & 22 & (7, 10) & 36 & (11, 15) \\
9 & (3, 6) & 23 & (7, 11) & 37 & (11, 16) \\
10 & (3, 7) & 24 & (7, 12) & 38 & (12, 15) \\
11 & (3, 8) & 25 & (8, 11) & 39 & (12, 16) \\
12 & (4, 7) & 26 & (8, 12) & 40 & (13, 14) \\
13 & (4, 8) & 27 & (9, 10) & 41 & (14, 15) \\
14 & (5, 6) & 28 & (9, 13) & 42 & (15, 16) \\
\bottomrule\end{longtable}
\section*{附录 B  核验与文件边界}
在项目根目录运行 python proofs/verify\_nonplanar\_rigidity.py，可从静态 GGB 的精确根式定义重建证书。加上 --export-full-construction 可重新导出完整代数区版本；不会覆盖原始 2\(\times\)2、静态参照或原联动文件。核验结果保存在 proofs/5x5 nonplanar rigidity certificate.json。

核验项目：36 个精确初始坐标、25 面的六距离、42\(\times\)36 的矩阵、-2\textsuperscript{36} 子式、R\(\bar v\)=0、\(\omega\)\textsuperscript{T}R=0、Q(\(\bar v\))=32，以及 20 步三球交点传播。数学证明本身依赖这些公开的有限代数恒等式，不依赖动画播放结果。

证书针对静态参照文件的 SHA-256 为：

dbd028ff36dd86b11ae8b5b5f8242ce0482de0893eae9c4cf01646518cd10262

\section*{参考文献}
[1] Robert Connelly and Walter Whiteley. Second-Order Rigidity and Prestress Stability for Tensegrity Frameworks. SIAM Journal on Discrete Mathematics, 9(3), 453-491, 1996. DOI: 10.1137/S0895480192229236. 本文相关条目：Theorem 4.3.1（印刷页 474）及 Corollary 5.2.2（印刷页 479）。

\href{https://pi.math.cornell.edu/~connelly/pdf/10.1137_S0895480192229236.pdf}{作者提供的论文 PDF（Cornell University）}
说明：文献提供一般刚性理论；本文的特定坐标、矩阵、整数权重、子式和 5\(\times\)5 推广由本项目计算并逐项验证。文献本身没有直接讨论这个折纸模型，也不能作为“所有非平面 5\(\times\)5 都不可折叠”的依据。

\end{document}
